\section*{Capítulo \ref{ch:b3:stem-cells} --- Células-tronco, regeneração e envelhecimento}

\begin{solution}{exo:b3:stem-cells:1}
\omterm{def:b3:stem-cells:stem}{Autorrenovação}: uma divisão que produz ao menos uma filha que ainda é
\omterm{def:b3:stem-cells:stem}{célula-tronco} --- as células \emph{Lgr5} da base da cripta. Potência: a
gama de destinos que uma célula ainda pode adotar --- a \omterm{def:b3:stem-cells:stem}{célula-tronco} da
cripta é multipotente, fazendo células absortivas, caliciformes,
enteroendócrinas, de Paneth e em tufo. \omterm{def:b3:stem-cells:stem}{Nicho}: as \omterm{prop:b3:stem-cells:crypt}{células de Paneth} e seus
sinais Wnt, Notch e EGF, que mantêm o caráter de tronco no lugar. \omterm{def:b3:stem-cells:stem}{Célula
amplificadora de trânsito}: uma filha comprometida que se divide mais
quatro ou cinco vezes na cripta antes de se diferenciar a caminho do topo
da vilosidade.
\end{solution}

\begin{solution}{exo:b3:stem-cells:2}
Células isoladas de medula produziram colônias esplênicas contendo
hemácias, granulócitos e plaquetas: uma célula, várias linhagens --- a
multipotência, e a resposta linear à dose mostrou que cada colônia vinha
de uma célula. A \omterm{def:b3:stem-cells:stem}{autorrenovação} exigiu um segundo camundongo: células de
uma colônia, injetadas em outro receptor irradiado, fizeram colônias
novas, de modo que a célula fundadora tinha produzido filhas com sua
própria capacidade.
\end{solution}

\begin{solution}{exo:b3:stem-cells:3}
\emph{Oct4}, \emph{Sox2}, \emph{Klf4}, \emph{c-Myc}. A \omterm{def:b3:chromatin-epigenetics:nucleosome}{cromatina} do
fibroblasto fechou os genes da \omterm{def:b3:stem-cells:pluripotency}{pluripotência} atrás de \omterm{def:b3:chromatin-epigenetics:nucleosome}{heterocromatina} e
\omterm{def:b3:chromatin-epigenetics:methylation}{metilação do DNA}, de modo que os fatores encontram poucos sítios
acessíveis a princípio e precisam recrutar enzimas de remodelamento ao
longo de muitos \omterm{def:b3:cell-cycle-apoptosis:cdk}{ciclos celulares}; a célula precisa ainda desligar seu
próprio programa, escapar da \omterm{def:b3:cell-cycle-apoptosis:other}{senescência} e passar por um estado
intermediário estocástico em que a maioria das células empaca ou morre.
Só a célula rara em que cada passo dá certo emerge pluripotente, semanas
depois.
\end{solution}

\begin{solution}{exo:b3:stem-cells:4}
A planária mantém células-tronco adultas pluripotentes, os \omterm{def:b3:stem-cells:regeneration}{neoblastos},
que migram para a ferida e constroem o que estiver faltando. A salamandra
não tem esse reservatório: células diferenciadas do coto se
desdiferenciam em \omterm{def:b3:stem-cells:stem}{progenitores} que, com células-tronco residentes do
tecido, formam um \omterm{def:b3:stem-cells:regeneration}{blastema} que roda de novo o programa do membro. Os
enxertos de tecidos isolados marcados feitos por Kragl mostraram que as
células do \omterm{def:b3:stem-cells:regeneration}{blastema} permanecem fiéis à sua origem --- músculo faz
músculo, cartilagem faz cartilagem ---, de modo que o \omterm{def:b3:stem-cells:regeneration}{blastema} é uma
mistura de \omterm{def:b3:stem-cells:stem}{progenitores} restritos a linhagens, e não uma massa
pluripotente.
\end{solution}

\begin{solution}{exo:b3:stem-cells:5}
$(11 - 5)/0.075 = 80$ duplicações. A partir de \qty{8}{kb}: $(8 - 5)/0.075 = 40$. Com \qty{60}{\%} de compensação a perda líquida é de
\qty{30}{bp}: $6000/30
= 200$.
\end{solution}

\begin{solution}{exo:b3:stem-cells:6}
Catorze divisões de células-tronco por dia rendem 28 filhas; 14 precisam
continuar células-tronco, de modo que 14 se comprometem --- uma filha
comprometida por \omterm{def:b3:stem-cells:stem}{célula-tronco} por dia. Cada uma se compromete a $2^{4} = 16$ células: $14\times 16 = 224$ células de vilosidade por dia, a ordem
das \num{300} observadas (uma quinta divisão de trânsito dá \num{448}). O
nível das células-tronco contribui com divisões equivalentes a catorze
células; o nível de trânsito, com as outras \num{200}.
\end{solution}

\begin{solution}{exo:b3:stem-cells:7}
$2\times 10^{11}/2^{20} = 1.9\times 10^{5}$ filhas comprometidas por dia;
com \num{10000} células-tronco, 19 por \omterm{def:b3:stem-cells:stem}{célula-tronco} por dia --- contra
uma divisão por ano. A contagem de vinte divisões é por linhagem, da
\omterm{def:b3:stem-cells:stem}{célula-tronco} à hemácia, mas os \omterm{def:b3:stem-cells:stem}{progenitores} intermediários não são de
uso único: os reservatórios de \omterm{def:b3:stem-cells:stem}{progenitores} multipotentes e comprometidos
se sustentam por semanas com suas próprias divisões, de modo que quase
toda a divisão acontece nos níveis de \omterm{def:b3:stem-cells:stem}{progenitor} e a \omterm{def:b3:stem-cells:stem}{célula-tronco} é
acionada raramente.
\end{solution}

\begin{solution}{exo:b3:stem-cells:8}
$\mu(50) = 10^{-3}\mathrm{e}^{1.74} = 5.7\times 10^{-3}$; $\mu(70) =
10^{-3}\mathrm{e}^{3.48} = 0.032$; $\mu(90) = 10^{-3}\mathrm{e}^{5.22} =
0.19$ por ano. Com $\mu_{0}/\gamma = 0.0115$: $S(70) = \exp[-0.0115
\times 31.5] = 0.70$ e $S(90) = \exp[-0.0115\times 184] = 0.12$. Mediana:
$30 + \ln(61)/0.087 = 77$. Reduzindo $\mu_{0}$ à metade:
$\ln(121)/0.087 = 55$, mediana 85 --- um tempo de duplicação, oito anos,
ganho. Reduzindo $\gamma$ à metade, a $0.0435$: $\ln(31)/0.0435 = 79$,
mediana 109: retardar o expoente ganha quatro vezes mais que reduzir a
constante à metade.
\end{solution}

\begin{solution}{exo:b3:stem-cells:9}
O terço restante precisa triplicar: $\log_{2}3 = 1.6$ divisão por célula.
\omterm{def:b3:stem-cells:regeneration}{Crescimento compensatório}, porque os lobos existentes aumentam pela
divisão de células no lugar; os lobos removidos, com seus ductos, vasos e
arquitetura, não são reconstruídos --- o fígado recupera sua massa e sua
função, mas não sua forma.
\end{solution}

\begin{solution}{exo:b3:stem-cells:10}
A cada evento o tamanho $k$ do clone vai a $k + 1$ ou a $k - 1$ com igual
probabilidade, de modo que seu tamanho esperado nunca muda; quando o
passeio termina, ele é $N$ (tomada) ou $0$ (perda), e a esperança $N\,P + 0 = k$ dá $P = k/N$. Uma célula única: $1/14$. Uma mutação neutra numa
\omterm{def:b3:stem-cells:stem}{célula-tronco} se fixa em sua cripta com probabilidade $1/14$ e, de outro
modo, é perdida em meses --- a deriva expurga treze de cada catorze
mutações.
\end{solution}

\begin{solution}{exo:b3:stem-cells:11}
O passeio agora é enviesado: o clone ganha um lugar com mais frequência
do que perde um, de modo que seu tamanho esperado cresce e a
probabilidade de tomada sobe rumo à certeza conforme $k$ cresce --- para
uma vantagem de \qty{10}{\%} num \omterm{def:b3:stem-cells:stem}{nicho} de catorze, cerca do triplo de
$1/14$ para uma única célula, e mais ainda quando o clone já ocupa alguns
lugares. Ao longo de décadas, as células-tronco com tais mutações
(\emph{DNMT3A}, \emph{TET2} na medula) tomam seus \omterm{def:b3:stem-cells:stem}{nichos} e, aos setenta,
um quinto das pessoas tem boa parte de seu sangue vindo de um clone. Não
é \omterm{def:b3:cancer-biology:hallmarks}{câncer}: as células ainda se diferenciam e obedecem à sua hierarquia.
Mas a mutação seguinte tem agora um clone de bilhões em que surgir, e o
risco de leucemia fica dez vezes maior.
\end{solution}

\begin{solution}{exo:b3:stem-cells:12}
Mediana $= 30 + \gamma^{-1}\ln(1 + \gamma\ln 2/\mu_{0})$: para $\mu_{0} =
10^{-2}$, $30 + \ln 7/0.087 = 52$; para $10^{-3}$, 77; para $10^{-4}$,
$30 +
\ln 604/0.087 = 104$. Na forma de Gompertz pura, cada corte de dez
vezes acrescenta os mesmos $\ln 10/\gamma = 26$ anos. Na prática os
retornos diminuem porque os ganhos históricos vieram de remover um termo
independente da idade --- infecção, parto, acidente, a constante $A$ de
Makeham em $\mu = A
+ \mu_{0}\mathrm{e}^{\gamma t}$ --- e, uma vez que $A$
é pequeno em comparação com a exponencial nas idades que importam,
cortá-lo mais quase nada muda; o $\mu_{0}$ intrínseco mudou muito menos.
Um corte de \qty{20}{\%} em $\gamma$ (para $0.070$) dá $30 + \ln 49/0.070 = 86$: nove anos, ganhos em toda idade, e não por adiar algumas
causas de morte.
\end{solution}

\begin{solution}{pb:b3:stem-cells:1}
\textbf{1.} $3.5\times 10^{11}$ por dia, $4\times 10^{6}$ por segundo.
\textbf{2.} $2^{20} \approx 10^{6}$ células por filha comprometida;
$3.5\times 10^{11}/10^{6} = 3.3\times 10^{5}$ filhas comprometidas por
dia.
\textbf{3.} $33$ por \omterm{def:b3:stem-cells:stem}{célula-tronco} por dia, \num{12000} por ano.
\textbf{4.} Os reservatórios de \omterm{def:b3:stem-cells:stem}{progenitores} precisam se sustentar: os
\omterm{def:b3:stem-cells:stem}{progenitores} multipotentes e comprometidos se dividem por semanas e
mantêm seus próprios números, de modo que quase toda divisão das vinte
acontece em níveis que se renovam, e a \omterm{def:b3:stem-cells:stem}{célula-tronco} contribui com uma
filha comprometida raramente --- da ordem de uma vez por ano.
\textbf{5.} $10 - 20\times 0.06 = \qty{8.8}{kb}$: bem acima de
$L_{\text{crit}}$, e a hemácia nunca mais se divide; a reserva só importa
para \omterm{def:b3:stem-cells:stem}{progenitores} que são reutilizados.
\textbf{6.} Perda líquida $0.3\times 60 = \qty{18}{bp}$; $6000/18 = 333$
divisões; a uma por ano, $333$ anos --- muito além de uma vida, de modo
que os telômeros não limitam a própria \omterm{def:b3:stem-cells:stem}{célula-tronco}; outros danos
limitam.
\textbf{7.} $14$ divisões por dia; $7$ repõem células-tronco perdidas e
$7$ se comprometem; $7\times 2^{4} = 112$ células de vilosidade por dia.
\textbf{8.} $N^{2} = 196$ eventos por célula a um por dia: cerca de
\qty{200}{d}, seis a sete meses --- o tempo em que as criptas do
confete ficaram de cor única.
\textbf{9.} $1/14 \approx \qty{7}{\%}$; $10^{7}/14 \approx 7\times 10^{5}$
criptas.
\textbf{10.} A deriva descarta treze de cada catorze mutações em poucos
meses. A \omterm{def:b3:stem-cells:stem}{divisão assimétrica} ao longo da vida manteria toda \omterm{def:b3:stem-cells:stem}{célula-tronco}
mutante para sempre, de modo que as mutações se acumulariam em toda
linhagem sem nunca serem expurgadas.
\textbf{11.} Ser tronco é uma posição no \omterm{def:b3:stem-cells:stem}{nicho} e um conjunto de sinais, e
não uma propriedade permanente de uma célula: qualquer célula que alcance
os sinais de Paneth pode assumir o papel.
\textbf{12.} Induza uma marcação multicolorida nas células-tronco num
dado momento e acompanhe as criptas. Deriva simétrica: as criptas ficam
de cor única em poucos meses. \omterm{def:b3:stem-cells:stem}{Divisão assimétrica}: cada cripta mantém
todas as suas cores indefinidamente, cada uma como uma fita fina
subindo a vilosidade.
\textbf{13.} $6000/60 = 100$ duplicações.
\textbf{14.} Cerca de $75$ restantes: a célula conta divisões, e não
tempo de calendário, já que a década no congelador nada custou.
\textbf{15.} $(8000 - 4000)/30 = 133$ anos depois dos vinte --- aos 153
anos. O comprimento médio dos telômeros dos leucócitos não limita, por si
só, a vida; os telômeros mais curtos, e as demais marcas, importam mais.
\textbf{16.} A \omterm{def:b3:cancer-biology:telomerase}{telomerase} remove uma barreira, a \omterm{thm:b3:stem-cells:hayflick}{senescência replicativa},
e por isso é encontrada em quase todo \omterm{def:b3:cancer-biology:hallmarks}{câncer}, que precisa de divisão
ilimitada; mas os fibroblastos que receberam \omterm{def:b3:cancer-biology:telomerase}{telomerase} mantiveram seus
pontos de checagem, a inibição por contato e a p53 intacta, de modo que a
divisão ilimitada sozinha não os tornou \omterm{def:b3:cancer-biology:hallmarks}{cânceres}. Necessária, não
suficiente.
\textbf{17.} $\mu(40) = 2.4\times 10^{-3}$, $\mu(60) = 0.014$, $\mu(80) =
0.077$, $\mu(100) = 0.44$ por ano. $S = \exp[-0.0115(\mathrm{e}^{\gamma
(t-30)} - 1)]$: $0.98$ aos 40, $0.87$ aos 60, $0.42$ aos 80 e $0.006$ aos
100.
\textbf{18.} Mediana $30 + \ln 61/0.087 = 77$. Dez por cento sobrevivem
em $\mathrm{e}^{\gamma(t-30)} = 1 + \gamma\ln 10/\mu_{0} = 201$: $t = 30 +
5.3/0.087 = 91$.
\textbf{19.} \qty{0.67}{mm/d}; $\log_{2}1000 = 10$ duplicações em 60
dias, uma a cada seis dias.
\textbf{20.} Cerca de $10^{8}$ divisões, portanto $0.1$ mutação
oncogênica por \omterm{def:b3:stem-cells:regeneration}{regeneração} --- uma em cada dez membros. Os axolotes quase
nunca desenvolvem \omterm{def:b3:cancer-biology:hallmarks}{tumores}: são de sangue frio e lentos, com supressão
tumoral robusta, e as células do \omterm{def:b3:stem-cells:regeneration}{blastema} permanecem restritas a
linhagens e se rediferenciam em vez de continuar se dividindo; o tecido
em \omterm{def:b3:stem-cells:regeneration}{regeneração} chega a suprimir \omterm{def:b3:cancer-biology:hallmarks}{tumores} induzidos. Um mamífero tem mais
células, temperatura e taxa metabólica mais altas e uma vida mais longa
ao longo da qual um clone escapado poderia crescer.
\textbf{21.} Desnerve o coto e implante uma esfera que libere a proteína
candidata, ou enxerte tecido condicionado por nervo: a \omterm{def:b3:stem-cells:regeneration}{regeneração}
restaurada significa um fator difusível. Candidatos identificados em
tritões e axolotes: a proteína de gradiente anterior nAG, a
neuregulina-1 e os FGFs.
\textbf{22.} A \omterm{prop:b3:stem-cells:crypt}{sinalização Wnt} é alta na cauda e baixa na cabeça, e a
ferida na extremidade anterior de um pedaço normalmente faz uma cabeça
porque o Wnt ali é baixo; silenciar um inibidor de Wnt eleva o Wnt em
toda parte, as duas feridas leem ``posterior'' e crescem duas caudas. O
pedaço conhece seu próprio eixo pelo gradiente residual de seu músculo, e
a ferida restabelece o gradiente em relação a ele.
\textbf{23.} O terço restante precisa triplicar: todo hepatócito se
divide uma vez (chegando a dois terços) e metade se divide de novo ---
uma média de $1.6$ divisão. Os lobos sobreviventes incham até a massa
original; os lobos faltantes e seus ductos e vasos não são reconstruídos.
\textbf{24.} Uma \omterm{def:b3:stem-cells:regeneration}{regeneração} que substituísse tecido gasto e danificado
retardaria o acúmulo que o expoente mede, baixando $\gamma$ e achatando a
curva --- os axolotes mostram pouca \omterm{def:b3:cell-cycle-apoptosis:other}{senescência}. O preço seria um termo
constante mais alto, vindo dos \omterm{def:b3:cancer-biology:hallmarks}{cânceres} num corpo quente, grande e
longevo.
\textbf{25.} Contagem de Hayflick de cerca de $100$ duplicações; uma
\omterm{def:b3:stem-cells:stem}{célula-tronco} auxiliada pela \omterm{def:b3:cancer-biology:telomerase}{telomerase} tem umas $330$ divisões, mais de
três séculos a uma por ano; expectativa mediana de vida de $77$ anos.
\end{solution}
